% Editable transcription; verify reused text against the original PDF.
% The abstract environment is supplied by main.tex.

Schank (1980) wrote an editorial for Intelligence on
``How much intelligence is there in artificial intelligence?''.
In this paper, we revisit this question.
We start with a short overview of modern AI and showcase some of
the AI breakthroughs in the four decades since Schank's paper.
We follow with a description of the main techniques these AI
breakthroughs were based upon, such as deep learning and reinforcement
learning; two techniques that have deep roots in psychology.
Next, we discuss how psychologically plausible AI is and could become
given the modern breakthroughs in AI's ability to learn.
We then access the main question of how intelligent AI systems
actually are.
For example, are there AI systems that can solve human intelligence
tests?
We conclude that Shank's observation, that intelligence is all about
generalization and that AI is not particularly good at this, has,
so far, withstood the test of time.
Finally, we consider what AI insights could mean for the study of
individual differences in intelligence.
We close with how AI can further Intelligence research and vice versa,
and look forward to fruitful interactions in the future.
